\chapter{Machines That Remain MSX and Amiga}
\label{ch:machines}

\section{The place of the original machines}

The MSX and the Amiga solved real-time picture and sound with small memories, narrow buses, and specialised peripheral chips. Those chips, and the way software talks to them, are the part worth keeping in register-transfer form. A student who can read a video display processor written as rules will understand a later display controller. A student who can read a shared sixteen-bit bus with several masters will understand a later DMA fabric.

BX-8 is the MSX-compatible member. Its slots, video display processor, programmable sound generator, keyboard matrix, and cartridge signalling follow the MSX map closely enough that a board can sit where an MSX sat: same kind of picture, same kind of sound, same habit of I/O ports. The processor is an RV32I hart with an eight-bit data path and the SVM-8 vector profile.

BX-16 is the Amiga-compatible member. Its data bus is sixteen bits wide. Memory is shared among the processor, a blitter-class mover, a video fetch, and an audio path, in the spirit of the Amiga chipset. The processor is an RV32I hart with a sixteen-bit data path and the SVM-16 vector profile. It can take the place of an Amiga in the same practical sense: a single FPGA image that runs a sixteen-bit desktop-class peripheral set and a vector unit sized for blits, mixing, and small transforms.

\section{What compatible means in RTL}

Compatibility, here, is a list you can tick on a board, not a promise that every original binary image will execute.

\begin{itemize}[leftmargin=1.5em]
  \item Address ranges and register layouts of the kept peripherals match the reference machines within documented differences.
  \item Bus width matches: eight bits on BX-8, sixteen bits on BX-16.
  \item Cycle behaviour of the peripherals is specified and tested against the reference documents you choose to follow.
  \item Software that is written for the BX hart, in RV32I and in the vector profile, drives those peripherals. This is the default elaboration, \texttt{host = RV32}.
  \item A thin compatibility layer, if you write one, translates or interprets guest code. It is software on the hart.
  \item A host elaboration attaches the same SVM to a softcore you already have: R800 on BX-8, 68000 or SuperH on BX-16, MIPS on BX-32. The adaptor recognises an escape and presents one OP-V parcel. It is not a second decoder hidden inside the RV32 pipeline. Each elaboration has one fetch path.
\end{itemize}

The original processors remain the reference for what those machines could do with a given number of cycles. In the RV32 elaboration they are not instantiated beside the hart. In a host elaboration the softcore is the scalar CPU, and the SVM is the vector unit beside it.

\section{Why the peripherals survive the change of processor}

A video display processor does not become simpler because the CPU is RISC-V. Sprites, tiles, blanking, and a shared memory port are the same problem on an MSX-compatible board as they were on the original. Writing them in Bluespec forces the arbitration into rules that the compiler schedules. That is the infrastructure BX-32 and BX-64 reuse when the picture is no longer a tilemap and the bus is no longer eight bits wide.

The same is true of the sixteen-bit machine. A blitter is a programmed DMA with a small arithmetic step. SVM-16 can perform the arithmetic in the vector unit, or the blitter can remain a peripheral master. Both organisations are legal on BX-16. The book recommends keeping the peripheral master for compatibility of timing, and using the vector unit for new software.

\section{The machine contract}

\begin{table}[htbp]
\centering
\caption{What each machine is for.}
\begin{tabular}{@{}llp{6.2cm}@{}}
\toprule
Machine & Can replace & Vector contract \\
\midrule
BX-8  & an MSX   & SVM-8, integer, EEW 8 \\
BX-16 & an Amiga & SVM-16, integer, EEW 8 and 16 \\
BX-32 & ---      & Zve32f, with FP32 \\
BX-64 & ---      & V or Zve64d, with FP64 \\
BX-128 & ---     & V or Zve64d, VLEN 128, ELEN 64 \\
BX-512 & ---     & V or Zve64d, VLEN 512, ELEN 64 \\
\bottomrule
\end{tabular}
\end{table}

BX-32 and BX-64 keep the organisational lessons of the Saturn, the Dreamcast, the PlayStation~2, and the Nintendo~64: more than one hart, local memory, and a vector unit wide enough for geometry. BX-128 and BX-512 keep that vector contract and pin the register length. RV128I, Chapter~\ref{ch:rv128}, is an integer base with XLEN 128. It has no row in this table and no vector profile. The default compatibility target from BX-32 upward is the RISC-V specification. A host elaboration can sit an SH-2, SH-4, R3000A, VR4300, or Emotion Engine softcore in front of the same SVM. That is an instruction port. It is not a claim that the FPGA image is the console, and it is not a cartridge port.
